% Background

\subsection{Three-State Totalistic Cellular Automata}
\label{sec:ca-rules}
We consider a two-dimensional grid of $N = N_r \times N_c$ agents, each in one of three
states: Buy~($s=0$), Hold~($s=1$), or Sell~($s=2$). At each discrete time step, every
interior agent updates its state deterministically based on the states of its four
cardinal (Von~Neumann) neighbors. Specifically, for an agent at position $(r,c)$ with
neighbor states $\{s_N, s_S, s_E, s_W\}$, the neighborhood sum is:
\begin{equation}
  c_{r,c} = s_N + s_S + s_E + s_W \in \{0, 1, \ldots, 8\}.
  \label{eq:neighborhood-sum}
\end{equation}
A \emph{totalistic rule} is a mapping $\mathcal{R}: \{0,1,\ldots,8\} \to \{0,1,2\}$ that
assigns a next state to each of the nine possible neighborhood sums. Since each of the
nine outputs can take three values, there are $3^9 = 19{,}683$ distinct rules. Each rule
is identified by its decimal index:
\begin{equation}
  \text{Rule index} = \sum_{j=0}^{8} a_j \cdot 3^j,
  \label{eq:rule-index}
\end{equation}
where $a_j = \mathcal{R}(j)$ is the next state when the neighborhood sum equals~$j$.

\subsection{Edge Agent Decision Model}
\label{sec:edge-agents}
Boundary cells serve as \emph{edge agents} that inject external information into the
grid. A fraction $f_e$ of boundary cells (selected uniformly at random with a fixed seed)
are designated as edge agents; the remaining boundary cells follow the same CA rule as
interior agents.

Each edge agent classifies the current price $P_\tau$ relative to an exponential moving
average (EMA) and its associated volatility band:
\begin{align}
  \hat{S}_\tau &= \hat{S}_{\tau-1} + \alpha\bigl(P_\tau - \hat{S}_{\tau-1}\bigr), \label{eq:ema} \\
  \hat{\sigma}_\tau &= \hat{\sigma}_{\tau-1} + \alpha\bigl(|P_\tau - \hat{S}_\tau| - \hat{\sigma}_{\tau-1}\bigr), \label{eq:ema-vol}
\end{align}
where $\alpha = 1/60$ corresponds to a 60-minute smoothing window. The edge agent state
is then:
\begin{equation}
  e_\tau = \begin{cases}
    0 \;\text{(Buy)}  & \text{if } P_\tau < \hat{S}_\tau - k\,\hat{\sigma}_\tau, \\
    1 \;\text{(Hold)} & \text{if } \hat{S}_\tau - k\,\hat{\sigma}_\tau \le P_\tau \le \hat{S}_\tau + k\,\hat{\sigma}_\tau, \\
    2 \;\text{(Sell)} & \text{if } P_\tau > \hat{S}_\tau + k\,\hat{\sigma}_\tau,
  \end{cases}
  \label{eq:edge-state}
\end{equation}
where $k$ controls the width of the hold band.

\subsection{Rule Taxonomy}
% AUTO-GENERATED by Table1-RuleTaxonomy.jl — do not edit by hand
\begin{table}[h]
\centering
\caption{Wolfram totalistic CA rule taxonomy. Base-3 encoding reads
right-to-left: digit $a_i$ gives the next state when the neighbourhood
sum index equals $i$.}
\label{tab:rule-taxonomy}
\begin{tabular}{ccll}
\toprule
Rule index & Base-3 encoding & Economic interpretation & Role \\
\midrule
6158 & $(022110002)_{3}$ & Imitation / herding & Homogeneous model; stable wave patterns \\
19280 & $(222110002)_{3}$ & Majority rule & 80\% of agents in crash model; promotes herding \\
19286 & $(222110022)_{3}$ & Minority / contrarian & 20\% of agents in crash model; destabilises grid \\
\bottomrule
\end{tabular}
\end{table}

Table~\ref{tab:rule-taxonomy} lists the three rules used in this study. Rule~19280
is a majority/herding rule: agents tend to adopt the dominant state in their
neighborhood, promoting coordination. Rule~19286 is a contrarian rule that differs
from 19280 only in its response to extreme sell neighborhoods
($c \in \{7,8\}$): when all neighbors are selling, a Rule~19286 agent \emph{sells}
(amplifying the cascade) rather than buying (damping it). This single-digit difference
in the base-3 encoding is the minimal perturbation that converts a stabilizing agent
into a destabilizing one.
